%
% $Id: $
%
%
% Compilar a .pdf con LaTeX (pdflatex)
% Es necesario instalar Beamer (paquete latex-beamer en Debian)
%

%
% Gr�ficos:
% Los gr�ficos pueden suministrarse en PNG, JPG, TIF, PDF, MPS
% Los EPS deben convertirse a PDF (usar epstopdf)
%

\documentclass{beamer}
\usetheme{Warsaw}
\usebackgroundtemplate{\includegraphics[width=\paperwidth]{format/libresoft-bg-soft.png}}
\usepackage[spanish]{babel}
\usepackage[latin1]{inputenc}
\usepackage{graphics}
\usepackage{amssymb} % Simbolos matematicos

%\definecolor{libresoftgreen}{RGB}{162,190,43}
%\definecolor{libresoftblue}{RGB}{0,98,143}

%\setbeamercolor{titlelike}{bg=libresoftgreen}

%% Metadatos del PDF.
\hypersetup{
  pdftitle={Desarrollo de aplicaciones web con Python y Django},
  pdfauthor={Jes�s M. Gonz�lez Barahona},
  pdfcreator={GSyC/LibreSoft, Universidad Rey Juan Carlos},
  pdfproducer=PDFLaTeX,
  pdfsubject={Taller con alumnos del IES La Arboleda},
}
%%

\begin{document}

\title{Desarrollo de aplicaciones web con Python y Django}
\subtitle{De cero a ciento y pico en poco tiempo}
\author{Jes�s M. Gonz�lez Barahona}
\institute{jgb@gsyc.es \\
\url{http://identi.ca/jgbarah} \url{http://twitter.com/jgbarah} \\
GSyC/LibreSoft, Universidad Rey Juan Carlos}

\date{Fuenlabrada, 20 de junio de 2011}

\frame{
\maketitle
\begin{center}
\includegraphics[width=6cm]{format/gsyc-urjc}
\end{center}
}


% Si el titulo o el autor se quieren acortar para los pies de p�gina
% se pueden redefinir aqu�:
%\title{Titulo corto}
%\author{Autores abreviado}


% LICENCIA DE REDISTRIBUCION DE LAS TRANSPAS
\frame{
~
\vspace{4cm}

\begin{flushright}
\copyright 2011 Jes�s M. Gonz�lez Barahona. \\

  Algunos derechos reservados.\\
  Este art�culo se distribuye bajo
  la licencia \\
  ``Reconocimiento-CompartirIgual 3.0 Espa�a'' de Creative Commons, \\
  disponible en \\
  \url{http://creativecommons.org/licenses/by-sa/3.0/es/deed.es}
\end{flushright}
}


%\section{}

%%---------------------------------------------------------------

\begin{frame}
\frametitle{�En qu� vamos a malgastar la pr�xima hora?}

\begin{itemize}
\item Haremos un servidor web sencillito
\item Hablaremos con �l
\item Lo complicaremos un poco
\item Esconderemos la complicaci�n, complicando m�s las cosas
\item Volveremos a hablarle
\item Si da tiempo, haremos Ajax y todo...
\end{itemize}

\end{frame}

%%---------------------------------------------------------------

\begin{frame}[fragile]
\frametitle{Un servidor web muuuuuuy sencillo}

\begin{itemize}
\item \url{servidor-http-simple.py}

\begin{verbatim}
import socket
mySocket = socket.socket(socket.AF_INET,
                         socket.SOCK_STREAM)
mySocket.bind(('localhost', 1234))
mySocket.listen(5)
while True:
    (recvSocket, address) = mySocket.accept()
    print recvSocket.recv(1024)
    recvSocket.send("HTTP/1.1 200 OK\r\n\r\n" +
      "<html><body><h1>Hello!</h1></body></html>" +
      "\r\n")
    recvSocket.close()
\end{verbatim}

\item Para ejecutarlo:

\begin{verbatim}
python servidor-http-simple.py
\end{verbatim}

\end{itemize}

\end{frame}

%%---------------------------------------------------------------

\begin{frame}[fragile]
\frametitle{�C�mo hablamos con �l?}

\begin{itemize}
\item Pues desde un navegador... (probamos con Firefox)
\item O desde un plugin espec�fico (probamos Web Developer de Firefox)
\item Pero tambi�n podemos ver qu� nos dice

\item O hablar HTTP ``a pelo'' (desde un telnet)
\begin{verbatim}
telnet localhost 1234
\end{verbatim}

\item O pedirle al amigo curl que hable por nosotros
\begin{verbatim}
curl http://localhost:1234/
\end{verbatim}

\item O al amigo wget
\begin{verbatim}
wget http://localhost:1234/
\end{verbatim}

\item Y podemos espiar cualquier conversaci�n con wireshark...
\end{itemize}

\end{frame}


%%---------------------------------------------------------------

\begin{frame}[fragile]
\frametitle{Un caza-robots}

\begin{itemize}
\item \url{servidor-http-random.py}

\begin{small}
\begin{verbatim}
import socket
import random
mySocket = socket.socket(socket.AF_INET,
                         socket.SOCK_STREAM)
mySocket.bind(('localhost', 1234))
mySocket.listen(5)
random.seed()
while True:
   (recvSocket, address) = mySocket.accept()
   print recvSocket.recv(1024)
   nextPage = str (random.randint (0,10000))
   nextUrl = "http://localhost:"+str(myPort)+"/"+nextPage
   htmlBody = "<p>Next page: <a href="' \
      + nextUrl + '">' + nextPage + "</a></p>"
   recvSocket.send("HTTP/1.1 200 OK \r\n\r\n" +
      "<html><body>" + htmlBody + "</body></html>" +
      "\r\n")
   recvSocket.close()
\end{verbatim}
\end{small}
\end{itemize}

\end{frame}


%%---------------------------------------------------------------

\begin{frame}[fragile]
\frametitle{Y ahora con Django 1.2.3}

\begin{verbatim}
export PATH=/usr/local/django/django/bin:$PATH
export PYTHONPATH=/usr/local/django:$PATH
cd myproject
python manage.py runserver 1234
\end{verbatim}

\begin{itemize}
\item Una aplicaci�n realmente simple: \url{myproject} \\
\url{http://localhost:1234/} \\
\url{http://localhost:1234/hello} \\
\url{http://localhost:1234/bye/Maria} \\
\url{http://localhost:1234/number/3}

\item Un gestor de contenidos bien simple: \url{cms_users}

\end{itemize}

\end{frame}


%%---------------------------------------------------------------

\begin{frame}[fragile]
\frametitle{Y el no va m�s: �un poco de Ajax!}

\begin{itemize}
\item Una p�gina que genera frases: \\
  \url{sentences_generator.html}
  \begin{itemize}
  \item JavaScript encastrado en el HTML
  \item Se ejecuta en el navegador al cargar la p�gina
  \end{itemize}
\item Una aplicacioncilla Ajax (bien simple): \\
  \url{words_provider} \\
  {\small \url{http://localhost:1234/apps/sentences_generator.html}}
  \begin{itemize}
  \item Sevidor Django que sirve:
    \begin{itemize}
    \item Una p�gina HTML con JavaScript encastrado
    \item Un canal con partes de frases
    \end{itemize}
  \item Se ejecuta el JavaScript en el navegador
  \item Pide partes de frases al servidor
  \item Muestra el resultado
  \end{itemize}

\end{itemize}

\end{frame}


%%---------------------------------------------------------------

\begin{frame}
\frametitle{Algunas URLs}

\begin{itemize}
\item Python: \url{http://python.org}
\item Django: \url{http://djangoproject.com}
\item Para aprender Python (por ejemplo): \\
  Dive into Python \\
  \url{http://diveintopython.org}
\item Para aprender Django (por ejemplo): \\
  The Django Book \\
  \url{http://djangobook.com/}
\item ``Tecnolog�as y Servicios para Aplicaciones en Internet'', \\
  asignatura con materiales libremente accesibles, \\
  su temario incluye el contenido de esta presentaci�n \\
\end{itemize}
  {\small \url{http://docencia.etsit.urjc.es/moodle/course/view.php?id=3}}

\end{frame}



\end{document}
